% Fragmento público generado desde propietarios íntegros.
% Residencia: 52.8.4.
% Fuente: ../incoming/radion_monografia_20260827/manuscrito/sections/09_sector_90120.tex:31-54
\paragraph{Extracción de la componente \(C^*/6\) del registro dodecafásico}

Los doce sectores se agrupan en seis parejas opuestas:
\begin{equation}
p_i=e_i+e_{i+6},\qquad
a_i=e_i-e_{i+6},\qquad i=1,\ldots,6.
\label{eq:sexadic-pairs}
\end{equation}
La reconstrucción tiene perfil \(1+5+6=12\). La apertura por pareja es
\begin{equation}
h=\frac{C^*}{6}.
\label{eq:h-C6}
\end{equation}
No es una normalización exterior ni la sexta parte de un blanco: es el índice
interno de la reconstrucción sexádica del registro \(\Rstate\).

Definimos
\begin{equation}
q_{h,-}=\exp\left[-\frac\pi{180}(A-h)\right],
\qquad
q_{h,+}=\exp\left[-\frac\pi{180}(A+h)\right].
\label{eq:q-h}
\end{equation}

